\section{Proofs for the second difference and the rounding error}\label{app:signs}
We keep the notation of Section~\ref{sec:shape}. This appendix proves
Lemma~\ref{lem:rounding} of Section~\ref{sec:shape-end} and the results of
Section~\ref{sec:convex} on the signs of the second difference.

\begin{proof}[Proof of Lemma~\ref{lem:rounding}]
We use the standard model of floating-point
arithmetic~\cite[Section~2.2]{higham2002}. If $a$ and $b$ are floating-point
numbers and $\circ$ is addition, multiplication or division, then the
computed value of $a\circ b$ is $(a\circ b)(1+\delta)$ with
$\lvert\delta\rvert\le u$, as long as the result does not overflow or
underflow. Underflow means that the exact result is nonzero and below the
smallest normal number $2^{-1022}$, and then the rounded result is at most
$2^{-1022}$. So a computed result that is at least $2^{-1000}$ did not
underflow, and by hypothesis this covers every operation in the recursion.
The integers in \eqref{eq:weights} are positive and at most $N(N+2)<2^{53}$,
so they are floating-point numbers, and each computed weight is the true
weight times some $1+\delta$.

We claim that $\hat f_t(b)=f_t(b)(1+\vartheta)$ with
$(1-u)^{3(t-1)}\le1+\vartheta\le(1+u)^{3(t-1)}$. We use induction on $t$,
and the case $t=1$ is clear. In the step from $t$ to $t+1$, each term on the
right of \eqref{eq:forward} is computed as $\hat f_t(b')$ times the true
weight times $(1+\delta_1)(1+\delta_2)$, and the sum of two terms gets one more
factor $1+\delta_3$. All numbers are positive. So $\hat f_{t+1}(b)$ lies
between $(1-u)^3$ and $(1+u)^3$ times the right side of \eqref{eq:forward}
with $\hat f_t$ in place of $f_t$, and the claim follows. Nothing overflows
on the way, because the true weights and the $f_t(b)$ are at most $1$ and
$(1+u)^{3N}<2$.

In the last division both numbers lie between $2^{-1000}$ and $2$. So the
quotient lies between $2^{-1001}$ and $2^{1001}$, it neither underflows nor
overflows, and the division adds one more factor $1+\delta$. Doubling is
exact. Hence $(C_N/E_N)/\hat R$ lies between
$(1-u)^{3(N-1)}(1+u)^{-3(N-1)-1}$ and $(1+u)^{3(N-1)}(1-u)^{-3(N-1)-1}$.
Since $1+u\le(1-u)^{-1}$, this gives both bounds with the exponent
$6(N-1)+1=6N-5$. For $N\le1686$ put $k=6N-5\le10111$. Then
$(1-u)^{-k}\le(1-ku)^{-1}\le1+2ku\le1+3\cdot10^{-12}$.
\end{proof}

\begin{proof}[Proof of Lemma~\ref{lem:interleave}]
We insert the numbers $\eta_i$ into $\zeta$ one at a time. If $\eta_i=0$, the
sign pattern does not change. If $\eta_i>0$, then $\zeta_i>0$ or
$\zeta_{i+1}>0$, so $\eta_i$ is placed next to an entry of its own sign, and
the pattern does not change. The case $\eta_i<0$ is the same. This proves the
first claim. Deleting entries from a list replaces its sign pattern by a
subsequence, and $\eta$ is obtained from $w$ by deleting the entries
$\zeta_i$. This proves the second claim.
\end{proof}

\begin{proof}[Proof of Lemma~\ref{lem:level-one}]
Put $\xi_t=P_t(1)/P_t(0)$. By \eqref{eq:P-rec}, $P_{t+1}(0)=(1-\varepsilon)P_t(0)$
and $P_{t+1}(1)=(1-\varepsilon-\frac ct)P_t(1)+\varepsilon P_t(0)$. So
\[
 \xi_{t+1}-\xi_t=\frac1{1-\varepsilon}\Bigl(\varepsilon-\frac{c\xi_t}t\Bigr),
 \qquad\xi_2=\frac{2\varepsilon}{1-\varepsilon}.
\]
We claim that $\xi_t<\varepsilon t/c$ for $t\ge2$. For $t=2$ this holds
because $c<1-\varepsilon$. If it holds for $t$, then
$0<\xi_{t+1}-\xi_t\le\frac\varepsilon{1-\varepsilon}<\frac\varepsilon c$,
and it holds for $t+1$. So $\xi_t$ is strictly increasing for $t\ge2$, and
$\xi_t<1$ whenever $t\le c/\varepsilon$.

Now \eqref{eq:D-rec} with $a=1$, and $D_t(0)=P_t(1)-2P_t(0)$, give for $t\ge2$
\[
 s_{t+1}=\bigl(1-r_t(2)\bigr)s_t+\sigma_t,\qquad
 \sigma_t=\Bigl(\varepsilon-\frac ct\Bigr)(\xi_t-2)\,P_t(0).
\]
Here $1-r_t(2)\ge1-\varepsilon-c=\varepsilon>0$. If $t<c/\varepsilon$, both
factors of $\sigma_t$ are negative and $\sigma_t>0$. If $t>c/\varepsilon$, the
first factor is positive, and the second is first negative and then positive,
because $\xi_t$ increases. So
$\operatorname{sp}(\sigma_2,\sigma_3,\dots)\preceq\Pi_3$. Let
$\Phi_t=\prod_{u=2}^{t-1}(1-r_u(2))>0$. Then
\[
 \frac{s_t}{\Phi_t}=s_2+\sum_{u=2}^{t-1}\frac{\sigma_u}{\Phi_{u+1}},\qquad
 s_2=1-3\varepsilon .
\]
If $\varepsilon<\frac13$, the terms $s_2,\sigma_2,\sigma_3,\dots$ are first
nonnegative, then nonpositive, then nonnegative. Their partial sums are
nonnegative and nondecreasing in the first stretch, nonincreasing in the
second, and nondecreasing in the third. So the sign pattern of the partial
sums is a subsequence of $\Pi_3$. If $\varepsilon\ge\frac13$, then
$c/\varepsilon\le1$, so the terms are first nonpositive and then nonnegative,
and the sign pattern of the partial sums is a subsequence of $(-,+)$. This
proves the first claim.

For the second claim we compute $P_3(0)=\frac12(1-\varepsilon)^2$ and
$P_3(1)=P_3(2)=\frac12-P_3(0)$, so $s_3=P_3(0)-P_3(1)=(1-\varepsilon)^2-\frac12$.
Suppose that $s_3=0$. Then $\varepsilon=1-2^{-1/2}$, $P_3(0)=P_3(1)=\frac14$
and $s_4=\sigma_3=(\varepsilon-\frac c3)(P_3(1)-2P_3(0))=-\frac{5\varepsilon-1}{12}$.
This is not $0$, because $\varepsilon\ne\frac15$.
\end{proof}

\begin{proof}[Proof of Theorem~\ref{thm:five-blocks}]
For $3\le t\le N$ let $Z_t$ be the list
\[
 Z_t=\bigl(s_N,s_{N-1},\dots,s_{t+1},\ s_t,D_t(2),D_t(3),\dots,D_t(t-2),s_t,\
 s_{t+1},\dots,s_{N-1},s_N\bigr).
\]
It has $2N-t-1$ entries. Its middle part is
$(D_t(1),D_t(2),\dots,D_t(t-1))$, because $D_t(t-1)=D_t(1)=s_t$ by symmetry.
The list $Z_3$ is $(s_N,\dots,s_4,s_3,s_3,s_4,\dots,s_N)$, and $Z_N$ is the
list in the statement.

We check that $Z_{t+1}$ is obtained from $Z_t$ as in
Lemma~\ref{lem:interleave}, for $3\le t\le N-1$. The first $N-t$ entries of
$Z_{t+1}$ are $s_N,\dots,s_{t+1}$, and they are copies of the first $N-t$
entries of $Z_t$. So $A_i=1$ and $B_i=0$ there. The next $t-2$ entries are
$D_{t+1}(b)$ for $2\le b\le t-1$. By \eqref{eq:D-rec},
\[
 D_{t+1}(b)=r_t(b-2)\,D_t(b-1)+\bigl(1-r_t(b+1)\bigr)D_t(b),
\]
and $D_t(b-1)$ and $D_t(b)$ are neighbors in the middle part of $Z_t$, at the
right places. The coefficients satisfy $r_t(b-2)\ge\varepsilon$ and
$1-r_t(b+1)\ge1-\varepsilon-c=\varepsilon$, since $2\le b\le t-1$. The last
$N-t$ entries of $Z_{t+1}$ are $s_{t+1},\dots,s_N$, and each is a copy of the
entry of $Z_t$ one place to its right. So $A_i=0$ and $B_i=1$ there.

By Lemma~\ref{lem:interleave},
$\operatorname{sp}(Z_N)\preceq\operatorname{sp}(Z_{N-1})\preceq\dots\preceq\operatorname{sp}(Z_3)$.
By Lemma~\ref{lem:level-one} the sign pattern $p$ of $(s_3,\dots,s_N)$ is a
subsequence of $\Pi_3$, so it is one of the lists $()$, $(+)$, $(-)$, $(+,-)$,
$(-,+)$ and $(+,-,+)$. The sign pattern of $Z_3$ is $p$ reversed, followed by
$p$, with equal neighbors merged. In the six cases this is $()$, $(+)$, $(-)$,
$(-,+,-)$, $(+,-,+)$ and $(+,-,+,-,+)$, and each is a subsequence of $\Pi_5$.

For $f_N$ we argue in the same way with $D^f_t(m)=f_t(m+1)-2f_t(m)+f_t(m-1)$
and the lists
\[
 Z^f_t=\bigl(s^f_N,\dots,s^f_{t+1},\ s^f_t,D^f_t(2),\dots,D^f_t(t-1),\ d_t,d_{t+1},\dots,d_{N-1}\bigr),
 \qquad2\le t\le N,
\]
where $s^f_t=D^f_t(1)$ and $d_u=(1-r_u(u+1))D^f_u(u)=(\varepsilon-\frac
cu)f_u(u-1)$. The entry $D^f_{t+1}(t)$ of $Z^f_{t+1}$ equals
$r_t(t-2)D^f_t(t-1)+d_t$ by \eqref{eq:D-rec}, and the other entries are as
before. The proof of Lemma~\ref{lem:level-one} applies to $s^f_t$ without
change, with $\xi_t=f_t(1)/f_t(0)$, $\xi_2=\frac\varepsilon{1-\varepsilon}$
and $s^f_2=1-3\varepsilon$. The numbers $d_u$ are negative for
$u<c/\varepsilon$ and positive for $u>c/\varepsilon$. So
$\operatorname{sp}(Z^f_2)$ is a subsequence of $(+,-,+)$ followed by $(-,+)$,
and $Z^f_N=(D^f_N(1),\dots,D^f_N(N-1))$. As before
$\operatorname{sp}(Z^f_N)\preceq\operatorname{sp}(Z^f_2)$, and a subsequence of $(+,-,+)$
followed by $(-,+)$ is a subsequence of $\Pi_5$ once equal neighbors are merged.
\end{proof}

\begin{proof}[Proof of Theorem~\ref{thm:no-holes}]
By symmetry $D_N(N-a_1)>0>D_N(N-a_0)$. Suppose that $D_N(a)\le0$ for some $a$
with $a_1<a<N-a_1$.

\emph{The case $D_N(a)<0$.} Then the entries of $Z_N$ at $a_0<a_1<a<N-a_1<N-a_0$
have the signs $-,+,-,+,-$. So $(-,+,-,+,-)\preceq\operatorname{sp}(Z_N)\preceq\Pi_5$
by Theorem~\ref{thm:five-blocks}. This is impossible, because a subsequence of
$\Pi_5$ with five terms is $\Pi_5$ itself.

\emph{The case $D_N(a)=0$.} We use the lists $Z_t$ of the proof of
Theorem~\ref{thm:five-blocks} and number the entries of $Z_t$ from $1$ to
$2N-t-1$. So $D_t(b)$ stands at place $N-t+b$ for $1\le b\le t-1$. For
$3\le t\le N-1$ let $W_t$ be the list $w$ of Lemma~\ref{lem:interleave} for
$\zeta=Z_t$ and $\eta=Z_{t+1}$. Then
$\operatorname{sp}(W_t)=\operatorname{sp}(Z_t)\preceq\Pi_5$. The entry of
$Z_{t+1}$ at place $k$ is a combination of the entries of $Z_t$ at the places
$k$ and $k+1$, with nonnegative coefficients. We call these two entries its
parents. For $2\le b\le t-1$ the parents of $D_{t+1}(b)$ are $D_t(b-1)$ and
$D_t(b)$, and both coefficients are positive.

For $4\le t\le N$ let $B_t$ be the set of integers $b$ with
$\max(2,a-N+t)\le b\le\min(a,t-2)$. It is not empty, because $2\le a\le N-2$.
We prove by downward induction on $t$ that there are places $p_1<p_2<p_4<p_5$
in $Z_t$ such that the entries there have the signs $-,+,+,-$, that
$D_t(b)=0$ for every $b\in B_t$, and that all the places $N-t+b$ with
$b\in B_t$ lie strictly between $p_2$ and $p_4$. For $t=N$ we have
$B_N=\{a\}$, and we take the places $a_0$, $a_1$, $N-a_1$ and $N-a_0$.

Let the claim hold for $t+1$, where $4\le t\le N-1$, with the places
$p_1'<p_2'<p_4'<p_5'$. Let $b\in B_{t+1}$. Then $D_{t+1}(b)=0$ is a
combination of its two parents with positive coefficients. So the parents
both vanish, or they have strictly opposite signs. In the second case one
parent is negative. In $W_t$ it stands next to $D_{t+1}(b)$, hence between the
positive entries of $Z_{t+1}$ at the places $p_2'$ and $p_4'$, and these stand
between the negative entries at $p_1'$ and $p_5'$. So
$(-,+,-,+,-)\preceq\operatorname{sp}(W_t)\preceq\Pi_5$, which is impossible.
Hence $D_t(b-1)=D_t(b)=0$ for every $b\in B_{t+1}$. Every element of $B_t$ is
of the form $b$ or $b-1$ with $b\in B_{t+1}$, so $D_t$ vanishes on $B_t$.

The entry of $Z_{t+1}$ at $p_1'$ is negative, so one of its parents is
negative. Let $p_1\in\{p_1',p_1'+1\}$ be its place in $Z_t$, and choose $p_2$,
$p_4$ and $p_5$ in the same way. Then $p_1\le p_1'+1\le p_2'\le p_2$, and
$p_1\ne p_2$ because the signs differ. In the same way $p_4<p_5$. Let $q^-$ and
$q^+$ be the smallest and the largest place $N-t-1+b$ with $b\in B_{t+1}$. The
places $N-t+b$ with $b\in B_t$ lie between $q^-$ and $q^++1$. Since
$p_2\le p_2'+1\le q^-$ and $p_4\ge p_4'\ge q^++1$, and the entries at $p_2$
and $p_4$ are not zero, these places lie strictly between $p_2$ and $p_4$.
This proves the claim for $t$.

The induction step with $t=4$ shows that both parents of $D_5(b)$ vanish for every
$b\in B_5\subseteq\{2,3\}$. The parents of $D_5(2)$ are $s_4$ and $D_4(2)$,
and those of $D_5(3)$ are $D_4(2)$ and $s_4$. So $s_4=D_4(2)=0$. By
\eqref{eq:D-rec} with $t=3$ and $a=2$, and since $D_3(1)=D_3(2)=s_3$, we have
$D_4(2)=r_3(0)s_3+\bigl(1-r_3(3)\bigr)s_3=2\varepsilon s_3$. So $s_3=0$ as
well, which contradicts Lemma~\ref{lem:level-one}.
\end{proof}
